\documentclass[10pt]{article}
\usepackage{geometry}
\geometry{a4paper} 
\usepackage{graphicx}
\usepackage{fancyhdr}
\pagestyle{fancy}
\setlength\headheight{26pt}
\rhead{\includegraphics[width=3cm]{cecierj.jpg}}
\usepackage{lipsum}

\usepackage{xcolor}

\begin{document}
\noindent \textbf{Fundação CECIERJ - Consórcio CEDERJ}\\
\noindent \textbf{Tecnologia em Sistemas de Computação - UFF.}\\
\noindent \textbf{Resolução da AD1 de Matemática para Computação.}\\ % Para a AD2, basta alterar o número da AD
\noindent \textbf{Aluno(a): } $<$preencher com o nome do aluno$>$ \\ %manter o símbolo \\ após inserir o nome
\noindent \textbf{Matrícula: } $<$preencher com a matrícula do aluno$>$ \\ %manter o símbolo \\ após inserir a matrícula
\noindent \textbf{Pólo: } $<$preencher com a matrícula do aluno$>$

\vspace{1.3cm}
\textbf{Questão 01}\\ %exemplo de questão sem item

\lipsum[3] %texto de exemplo. Deletar antes de usar

\vspace{1.5em} % sempre que usar um ambiente para fórmulas fora de um texto use esse espaçamento
{\large $f(x) = \frac{x^2 - 4}{x - 2} = \frac{(x - 2)\times(x + 2)}{x - 2} = x + 2$ }%ambiente para fórmulas junto com texto

\vspace{1cm} %usar sempre o espaçamento de 1cm antes de começar uma nova questão
\textbf{Questão 02}\\ %exemplo de questão com item
\begin{itemize}
	\item[Item a] \lipsum[2] %sempre crie um novo item com o formato \item[Item x], sendo x a letra do item em questão.
\end{itemize}

\vspace{1cm}
\textbf{Questão 03}\\ %exemplo de questão sem item

\lipsum[3] %texto de exemplo. Deletar antes de usar

\end{document}